% Options for packages loaded elsewhere
\PassOptionsToPackage{unicode}{hyperref}
\PassOptionsToPackage{hyphens}{url}
\documentclass[
]{article}
\usepackage{xcolor}
\usepackage[margin=1in]{geometry}
\usepackage{amsmath,amssymb}
\setcounter{secnumdepth}{-\maxdimen} % remove section numbering
\usepackage{iftex}
\ifPDFTeX
  \usepackage[T1]{fontenc}
  \usepackage[utf8]{inputenc}
  \usepackage{textcomp} % provide euro and other symbols
\else % if luatex or xetex
  \usepackage{unicode-math} % this also loads fontspec
  \defaultfontfeatures{Scale=MatchLowercase}
  \defaultfontfeatures[\rmfamily]{Ligatures=TeX,Scale=1}
\fi
\usepackage{lmodern}
\ifPDFTeX\else
  % xetex/luatex font selection
\fi
% Use upquote if available, for straight quotes in verbatim environments
\IfFileExists{upquote.sty}{\usepackage{upquote}}{}
\IfFileExists{microtype.sty}{% use microtype if available
  \usepackage[]{microtype}
  \UseMicrotypeSet[protrusion]{basicmath} % disable protrusion for tt fonts
}{}
\makeatletter
\@ifundefined{KOMAClassName}{% if non-KOMA class
  \IfFileExists{parskip.sty}{%
    \usepackage{parskip}
  }{% else
    \setlength{\parindent}{0pt}
    \setlength{\parskip}{6pt plus 2pt minus 1pt}}
}{% if KOMA class
  \KOMAoptions{parskip=half}}
\makeatother
\usepackage{color}
\usepackage{fancyvrb}
\newcommand{\VerbBar}{|}
\newcommand{\VERB}{\Verb[commandchars=\\\{\}]}
\DefineVerbatimEnvironment{Highlighting}{Verbatim}{commandchars=\\\{\}}
% Add ',fontsize=\small' for more characters per line
\usepackage{framed}
\definecolor{shadecolor}{RGB}{248,248,248}
\newenvironment{Shaded}{\begin{snugshade}}{\end{snugshade}}
\newcommand{\AlertTok}[1]{\textcolor[rgb]{0.94,0.16,0.16}{#1}}
\newcommand{\AnnotationTok}[1]{\textcolor[rgb]{0.56,0.35,0.01}{\textbf{\textit{#1}}}}
\newcommand{\AttributeTok}[1]{\textcolor[rgb]{0.13,0.29,0.53}{#1}}
\newcommand{\BaseNTok}[1]{\textcolor[rgb]{0.00,0.00,0.81}{#1}}
\newcommand{\BuiltInTok}[1]{#1}
\newcommand{\CharTok}[1]{\textcolor[rgb]{0.31,0.60,0.02}{#1}}
\newcommand{\CommentTok}[1]{\textcolor[rgb]{0.56,0.35,0.01}{\textit{#1}}}
\newcommand{\CommentVarTok}[1]{\textcolor[rgb]{0.56,0.35,0.01}{\textbf{\textit{#1}}}}
\newcommand{\ConstantTok}[1]{\textcolor[rgb]{0.56,0.35,0.01}{#1}}
\newcommand{\ControlFlowTok}[1]{\textcolor[rgb]{0.13,0.29,0.53}{\textbf{#1}}}
\newcommand{\DataTypeTok}[1]{\textcolor[rgb]{0.13,0.29,0.53}{#1}}
\newcommand{\DecValTok}[1]{\textcolor[rgb]{0.00,0.00,0.81}{#1}}
\newcommand{\DocumentationTok}[1]{\textcolor[rgb]{0.56,0.35,0.01}{\textbf{\textit{#1}}}}
\newcommand{\ErrorTok}[1]{\textcolor[rgb]{0.64,0.00,0.00}{\textbf{#1}}}
\newcommand{\ExtensionTok}[1]{#1}
\newcommand{\FloatTok}[1]{\textcolor[rgb]{0.00,0.00,0.81}{#1}}
\newcommand{\FunctionTok}[1]{\textcolor[rgb]{0.13,0.29,0.53}{\textbf{#1}}}
\newcommand{\ImportTok}[1]{#1}
\newcommand{\InformationTok}[1]{\textcolor[rgb]{0.56,0.35,0.01}{\textbf{\textit{#1}}}}
\newcommand{\KeywordTok}[1]{\textcolor[rgb]{0.13,0.29,0.53}{\textbf{#1}}}
\newcommand{\NormalTok}[1]{#1}
\newcommand{\OperatorTok}[1]{\textcolor[rgb]{0.81,0.36,0.00}{\textbf{#1}}}
\newcommand{\OtherTok}[1]{\textcolor[rgb]{0.56,0.35,0.01}{#1}}
\newcommand{\PreprocessorTok}[1]{\textcolor[rgb]{0.56,0.35,0.01}{\textit{#1}}}
\newcommand{\RegionMarkerTok}[1]{#1}
\newcommand{\SpecialCharTok}[1]{\textcolor[rgb]{0.81,0.36,0.00}{\textbf{#1}}}
\newcommand{\SpecialStringTok}[1]{\textcolor[rgb]{0.31,0.60,0.02}{#1}}
\newcommand{\StringTok}[1]{\textcolor[rgb]{0.31,0.60,0.02}{#1}}
\newcommand{\VariableTok}[1]{\textcolor[rgb]{0.00,0.00,0.00}{#1}}
\newcommand{\VerbatimStringTok}[1]{\textcolor[rgb]{0.31,0.60,0.02}{#1}}
\newcommand{\WarningTok}[1]{\textcolor[rgb]{0.56,0.35,0.01}{\textbf{\textit{#1}}}}
\usepackage{graphicx}
\makeatletter
\newsavebox\pandoc@box
\newcommand*\pandocbounded[1]{% scales image to fit in text height/width
  \sbox\pandoc@box{#1}%
  \Gscale@div\@tempa{\textheight}{\dimexpr\ht\pandoc@box+\dp\pandoc@box\relax}%
  \Gscale@div\@tempb{\linewidth}{\wd\pandoc@box}%
  \ifdim\@tempb\p@<\@tempa\p@\let\@tempa\@tempb\fi% select the smaller of both
  \ifdim\@tempa\p@<\p@\scalebox{\@tempa}{\usebox\pandoc@box}%
  \else\usebox{\pandoc@box}%
  \fi%
}
% Set default figure placement to htbp
\def\fps@figure{htbp}
\makeatother
\setlength{\emergencystretch}{3em} % prevent overfull lines
\providecommand{\tightlist}{%
  \setlength{\itemsep}{0pt}\setlength{\parskip}{0pt}}
\usepackage[table]{xcolor}
\usepackage{multirow}
\usepackage{textgreek}
\usepackage{setspace}
\onehalfspacing
\usepackage{array}
\usepackage{colortbl}
\usepackage[pdfpagemode=UseNone,bookmarks=false,bookmarksopen=false,hidelinks]{hyperref}
\hypersetup{pdfmenubar=false,pdftoolbar=false,pdfstartview={XYZ null null 0.9}}
\usepackage{titling}
\pretitle{\begin{center}\LARGE\bfseries}
\posttitle{\par\vspace{1em} {\large Prepared by Mr. Patrick Mwaura (DSAS) \newline For \newline Mrs. Stella Nyambariga, Head of HR }\par\end{center}\vskip 7em}
\preauthor{}\postauthor{}
\predate{\begin{center}\large}\postdate{\par\end{center}}
\usepackage{bookmark}
\IfFileExists{xurl.sty}{\usepackage{xurl}}{} % add URL line breaks if available
\urlstyle{same}
\hypersetup{
  pdftitle={KIPRE STAFF TRAINING NEEDS ASSESSMENT (TNA) REPORT},
  hidelinks,
  pdfcreator={LaTeX via pandoc}}

\title{KIPRE STAFF TRAINING NEEDS ASSESSMENT (TNA) REPORT}
\author{}
\date{\vspace{-2.5em}2026-01-31}

\begin{document}
\maketitle

\definecolor{softred}{RGB}{255,200,200}
\definecolor{softyellow}{RGB}{255,255,200}
\definecolor{softgreen}{RGB}{200,255,200}
\definecolor{factorgray}{RGB}{220,220,220}

\section*{Executive Summary}

\newpage

\subsection{Background}\label{background}

The primary purpose of this Training Needs Analysis (TNA) is to
systematically assess institutional competency gaps and evaluate
workforce preparedness. This evidence-based approach is designed to
inform training priorities and guide strategic human capital development
across the KIPRE. The assessment was conducted using a multi-dimensional
questionnaire structured into specific modules to capture a holistic
view of the workforce. These sections included:

\begin{itemize}
\tightlist
\item
  Basic Information (Directorate, Age, and Sex of the employee)
\item
  Job Role and Responsibilities
\item
  Self Performances Assessing Questions
\item
  Skills and Copetency Assessment
\item
  Training History
\item
  Identification of Training Needs
\item
  Training Delivery Preferences
\item
  Organizational Support
\item
  Future Skills and Emerging Needs
\item
  Additional Comments
\end{itemize}

By triangulating data from these domains, the analysis identifies not
only current skill deficiencies but also the organizational support
structures and delivery formats required to bridge them effectively.

\subsection{Methodology}\label{methodology}

The assessment was digitized via ODK Central and deployed from
\(20^{th}-28^{th}\) April of 2026. To ensure we captured a
representative voice, we used a multi-channel approach---utilizing
official email and professional WhatsApp groups. Participation was
voluntary, positioned as a collaborative effort to co-design the
organization's future training roadmap. Data Processing \& Analysis Once
the collection window closed, raw data was ingested into R (v4.5.2). The
analytical pipeline focused on cleaning, recoding qualitative responses
for statistical depth, and aggregating the results to provide a granular
view of needs across all Directorates and Divisions.

\subsection{Workforce profile}\label{workforce-profile}

A total of 80 employees, representing 42\% of the 189 staff members
currently at KIPRE, participated in this assessment. This respondent
base spans 23 distinct divisions across seven primary Directorates,
ensuring the training roadmap reflects both the core scientific mandate
and the essential administrative backbone of the institute. The
distribution reveals a workforce concentrated within key technical hubs,
most notably the Animal Sciences Welfare and Ethics and Research and
Product Development Directorates. In the research domain, the Kenya
Snakebite Research and Intervention Division and the Welfare and Ethics
Division emerged as the most represented units. This engagement from
frontline technical staff---supported by a demographic where 51.2\% of
respondents are aged between 25 and 44---provides a robust foundation
for identifying gaps in laboratory biosafety and modern research
protocols. While the gender distribution is relatively balanced (45\%
female, 55\% male), distinct generational trends are evident: the female
cadre is most concentrated in the 35--44 age bracket (36.1\%), whereas
the male workforce peaks in the 45--54 bracket (34.1\%). Administrative
and strategic functions are equally well-mapped, with consistent
representation from Corporate Services and the Office of the Director
General. Notably, only 10\% of respondents are aged 55 or older,
indicating a predominantly mid-career workforce with significant
long-term growth potential. This demographic profile highlights a
strategic opportunity to implement cross-cutting training programmes
that standardise administrative excellence and ensure the institution is
prepared for future succession and evolving digital demands.

\begin{Shaded}
\begin{Highlighting}[]
\FunctionTok{library}\NormalTok{(tidyverse)}
\FunctionTok{library}\NormalTok{(janitor)}
\end{Highlighting}
\end{Shaded}

\begin{verbatim}
## 
## Attaching package: 'janitor'
\end{verbatim}

\begin{verbatim}
## The following objects are masked from 'package:stats':
## 
##     chisq.test, fisher.test
\end{verbatim}

\begin{Shaded}
\begin{Highlighting}[]
\CommentTok{\# 1. Summary by Directorate and Division}
\NormalTok{org\_summary }\OtherTok{\textless{}{-}}\NormalTok{ data }\SpecialCharTok{\%\textgreater{}\%}
  \FunctionTok{count}\NormalTok{(section\_a.directorate, section\_a.division\_section, }\AttributeTok{name =} \StringTok{"Staff\_Count"}\NormalTok{) }\SpecialCharTok{\%\textgreater{}\%}
  \FunctionTok{mutate}\NormalTok{(}
    \AttributeTok{Directorate =} \FunctionTok{str\_to\_title}\NormalTok{(}\FunctionTok{str\_replace\_all}\NormalTok{(section\_a.directorate, }\StringTok{"\_"}\NormalTok{, }\StringTok{" "}\NormalTok{)),}
    \AttributeTok{Division =} \FunctionTok{str\_to\_title}\NormalTok{(}\FunctionTok{str\_replace\_all}\NormalTok{(section\_a.division\_section, }\StringTok{"\_"}\NormalTok{, }\StringTok{" "}\NormalTok{))}
\NormalTok{  ) }\SpecialCharTok{\%\textgreater{}\%}
  \FunctionTok{select}\NormalTok{(Directorate, Division, Staff\_Count)}

\CommentTok{\# 2. Gender and Age Distribution}
\NormalTok{demographic\_summary }\OtherTok{\textless{}{-}}\NormalTok{ data }\SpecialCharTok{\%\textgreater{}\%}
  \FunctionTok{tabyl}\NormalTok{(section\_a.gender, section\_a.age\_bracket) }\SpecialCharTok{\%\textgreater{}\%}
  \FunctionTok{adorn\_totals}\NormalTok{(}\FunctionTok{c}\NormalTok{(}\StringTok{"row"}\NormalTok{, }\StringTok{"col"}\NormalTok{)) }\SpecialCharTok{\%\textgreater{}\%}
  \FunctionTok{adorn\_percentages}\NormalTok{(}\StringTok{"row"}\NormalTok{) }\SpecialCharTok{\%\textgreater{}\%}
  \FunctionTok{adorn\_pct\_formatting}\NormalTok{(}\AttributeTok{digits =} \DecValTok{1}\NormalTok{) }\SpecialCharTok{\%\textgreater{}\%}
  \FunctionTok{adorn\_ns}\NormalTok{() }\CommentTok{\# This shows both Count and Percentage (\%)}

\CommentTok{\# 3. Job Grade and Service Length}
\NormalTok{grade\_summary }\OtherTok{\textless{}{-}}\NormalTok{ data }\SpecialCharTok{\%\textgreater{}\%}
  \FunctionTok{count}\NormalTok{(section\_a.job\_grade, section\_a.length\_service) }\SpecialCharTok{\%\textgreater{}\%}
  \FunctionTok{pivot\_wider}\NormalTok{(}\AttributeTok{names\_from =}\NormalTok{ section\_a.length\_service, }\AttributeTok{values\_from =}\NormalTok{ n, }\AttributeTok{values\_fill =} \DecValTok{0}\NormalTok{)}

\CommentTok{\# Print results to console}
\FunctionTok{print}\NormalTok{(}\StringTok{"{-}{-}{-} Staff Count by Division {-}{-}{-}"}\NormalTok{)}
\end{Highlighting}
\end{Shaded}

\begin{verbatim}
## [1] "--- Staff Count by Division ---"
\end{verbatim}

\begin{Shaded}
\begin{Highlighting}[]
\FunctionTok{print}\NormalTok{(org\_summary)}
\end{Highlighting}
\end{Shaded}

\begin{verbatim}
## # A tibble: 23 x 3
##    Directorate                                              Division Staff_Count
##    <chr>                                                    <chr>          <int>
##  1 Animal Sciences Welfare And Ethics Directorate           Kenya S~           1
##  2 Animal Sciences Welfare And Ethics Directorate           Laborat~           1
##  3 Animal Sciences Welfare And Ethics Directorate           Nationa~           4
##  4 Animal Sciences Welfare And Ethics Directorate           Veterin~           9
##  5 Animal Sciences Welfare And Ethics Directorate           Welfare~          11
##  6 Capacity Building Partnership And Grant Management Dire~ Human R~           1
##  7 Capacity Building Partnership And Grant Management Dire~ Partner~           3
##  8 Corporate Services Directorate                           Corpora~           4
##  9 Corporate Services Directorate                           Finance~           4
## 10 Corporate Services Directorate                           Human R~           5
## # i 13 more rows
\end{verbatim}

\begin{Shaded}
\begin{Highlighting}[]
\FunctionTok{print}\NormalTok{(}\StringTok{"{-}{-}{-} Gender by Age Bracket (\%) {-}{-}{-}"}\NormalTok{)}
\end{Highlighting}
\end{Shaded}

\begin{verbatim}
## [1] "--- Gender by Age Bracket (%) ---"
\end{verbatim}

\begin{Shaded}
\begin{Highlighting}[]
\FunctionTok{print}\NormalTok{(demographic\_summary)}
\end{Highlighting}
\end{Shaded}

\begin{verbatim}
##  section_a.gender 25_34_years 35_44_years 45_54_years 55+_years under_25_years
##            female  22.2%  (8)  36.1% (13)  27.8% (10)  5.6% (2)       8.3% (3)
##              male  29.5% (13)  15.9%  (7)  34.1% (15) 13.6% (6)       6.8% (3)
##             Total  26.2% (21)  25.0% (20)  31.2% (25) 10.0% (8)       7.5% (6)
##        Total
##  100.0% (36)
##  100.0% (44)
##  100.0% (80)
\end{verbatim}

\begin{Shaded}
\begin{Highlighting}[]
\FunctionTok{print}\NormalTok{(}\StringTok{"{-}{-}{-} Job Grade by Length of Service {-}{-}{-}"}\NormalTok{)}
\end{Highlighting}
\end{Shaded}

\begin{verbatim}
## [1] "--- Job Grade by Length of Service ---"
\end{verbatim}

\begin{Shaded}
\begin{Highlighting}[]
\FunctionTok{print}\NormalTok{(grade\_summary)}
\end{Highlighting}
\end{Shaded}

\begin{verbatim}
## # A tibble: 11 x 6
##    section_a.job_grade `1_3_years` `16+_years` `4_7_years` `8_15_years`
##    <chr>                     <int>       <int>       <int>        <int>
##  1 nm1                           1           0           0            0
##  2 nm10                          1           5           2            4
##  3 nm11                          0           8           1            1
##  4 nm12                          5           2           1            1
##  5 nm13                          2           1           0            1
##  6 nm2                           1           0           0            0
##  7 nm5                           0           3           0            1
##  8 nm6                           0           2           0            0
##  9 nm7                           3           3           2            2
## 10 nm8                           3           4           9            3
## 11 nm9                           0           1           0            1
## # i 1 more variable: less_than_1_year <int>
\end{verbatim}

\subsubsection{1. Which critical competency gaps exist across
directorates, divisions, and job grades, and which areas require
immediate training
intervention?}\label{which-critical-competency-gaps-exist-across-directorates-divisions-and-job-grades-and-which-areas-require-immediate-training-intervention}

This question helps management identify

High-priority skill deficiencies Departments with the greatest capacity
gaps Staff categories requiring urgent support

Likely outputs

Competency gap heatmaps Directorate ranking tables Priority training
matrix Skill deficiency frequencies

Likely recommendations

Targeted capacity-building programmes Directorate-specific training
plans Prioritised budget allocation

\subsection{2. To what extent are employees adequately prepared to meet
emerging organisational, technological, and policy-related
demands?}\label{to-what-extent-are-employees-adequately-prepared-to-meet-emerging-organisational-technological-and-policy-related-demands}

\subsection{Key recommendations}\label{key-recommendations}

Priority actions

\subsection{Conclusion}\label{conclusion}

Strategic implications

\subsection{7. Examplary visualisations}\label{examplary-visualisations}

\end{document}
